\documentstyle[% 


righttag%


,ctagsplt%


,11pt% 


,amssymb% 


,russian%               remove in Eng. 


,izvru%                 Set izven% in Eng. 


% ,izvru-ks             for brief communications 


]{amsart} 


 


%       Math definitions 


 


\newcommand{\col}{\operatorname{col}} 


\newcommand{\vraisup}{\operatornamewithlimits{vrai\,sup}} 


\newcommand{\vraimax}{\operatornamewithlimits{vrai\,max}} 


\renewcommand{\baselinestretch}{1.13}





%%         Use for single entries 


%% \theoremstyle{plain} 


%% \theorembodyfont{\it} 


%% \newtheorem{thmnonum}{\hskip\parindent Теорема} 


%% \renewcommand{\thethmnonum}{} 


 


\theoremstyle{definition} 


\theorembodyfont{\rm} 


\newtheorem{notenonum}{\hskip\parindent Замечание} 


\renewcommand{\thenotenonum}{} 


 


%\hoffset=-15mm                  For Eng. 


 


\setcounter{page}{25} 


\god{1997} 


\nomer{6~(421)} %                Remove in Eng. 


%\nomer{Vol.\,41, No.\,6}        Set in Eng. 


% \str{75-81}                    Set in Eng. 


\udk{517.929} 


\author{\it А.Н.\,Румянцев} 


\title{Конструктивное исследование\\


дифференциальных уравнений\\


с произвольным отклонением аргумента}


 


%\thanks{Работа выполнена при поддержке РФФИ, грант \No. 94-01-00183.} 


\date{} 


% \pagestyle{myheadings}         Set in Eng. 


 


\pagestyle{plain} %              Renove in Eng. 


\begin{document} 


 


% \thispagestyle{plain}          Set in Eng. 


 


\maketitle 


 


% Body text 


 


Статья продолжает цикл работ по развитию конструктивного метода в теории


функциональ\-но-дифференциальных уравнений


\cite{Az}--\cite{Mak2}. Объектом настоящего


исследования является линейное дифференциальное уравнение


с произвольным отклонением аргумента.





Будет уместным указать на ряд работ зарубежных авторов, относящихся к этому


направлению. Работы


\cite{Kau1}, \cite{Kau2} посвящены построению конструктивных методов


исследования обыкновенных дифференциальных и интегральных уравнений. В


\cite{Moo} рассматриваются интервальные методы анализа операторных уравнений. В


работе \cite{Plum} предлагаются конструктивные методы


исследования разрешимости нелинейных краевых задач для обыкновенных


дифференциальных уравнений.





В вводной части работы дается краткий обзор необходимых


сведений из теории функциональ\-но-дифференциальных уравнений.


Далее описывается обоснование конструктивного, ориентированного на


применение компьютера, метода исследования существования решения для


дифференциального уравнения с произвольным отклонением аргумента.


В заключение приводится иллюстрирующий пример.





\section{Обозначения}





Всюду ниже: $R$ --- пространство вещественных чисел с нормой $|\cdot|$; $L$


--- пространство суммируемых функций $z:[0,T]\to R$;


$\|z\|_L=\int_0^T |z(s)|ds$; $L_\infty$ --- пространство функций


$z:[0,T] \to R$;


$\|z\|_{L_\infty}=\vraisup\limits_{t \in [0,T]} |z(t)|$;


$D$ --- пространство абсолютно


непрерывных функций $x:[0,T] \to R$;


$\|x\|_D=|x(0)| + \|\Dot x \|_L$; $S_h:D \to L$ -- оператор внутренней


суперпозиции:


$$


(S_h)(x)=


\begin{cases}


x[h(t)], & \text{ если } h(t) \in [0,T]; \\


0,       & \text{ если } h(t) \notin [0,T],


\end{cases}


$$


функция $\phi^h$ определяется равенством


$$


(\phi^h)(t)=


\begin{cases}


0,           & \text{ если } h(t) \in [0,T]; \\


\phi [h(t)], & \text{ если } h(t) \notin [0,T].


\end{cases}


$$





\section{Введение}





Приведем необходимые сведения из теории линейных


функционально-дифференциальных уравнений \cite{Az}.





Рассматривается уравнение


%


\begin{equation}


({\cal L}x)(t) =  f(t), \quad t \in [0,T];


\end{equation}


%


$f \in L$; ${\cal L}: D \to L$ --- линейный ограниченный оператор.


Для всякого $x \in D$ имеет место представление 


$x(t) = \int_0^t \Dot x(s) ds  + x(0)$,


в силу которого уравнение (1) можно записать в виде


%


\begin{equation}


({\cal L}x)(t) = (Q \Dot x)(t) + A(t)x(0) = f(t), \quad


t \in [0,T],


\end{equation}


%


где оператор $Q:L \to L$; $ (Qz)(t) = {\cal L} \Bigl(


\int_0^{(\cdot ) }\Dot x(s) ds \Bigr)$, называется главной


частью оператора ${\cal L}$; $A(t)=({\cal L} e)(t)$; $e(t)=1$,


$t\in [0,t]$. Для широкого класса функционально-дифференциальных уравнений


оператор $Q$ фредгольмов и имеет представление $Q=I-K$, где $K:L \to L$ ---


интегральный вполне непрерывный оператор: 


$(Kz)(t) =\int_0^T k(t,s)z(s)ds$; $I$ --- тождественный оператор. Приведем


примеры таких уравнений.





ПРИМЕР 1. Обыкновенное дифференциальное уравнение


$$


({\cal L}x)(t) \equiv \Dot x(t) + p(t)x(t) = f(t),


\quad p \in L, \quad t \in [0,T].


$$


В этом случае $(Kz)(t) = - \int_0^t p(t)z(s) ds$.





ПРИМЕР 2. Дифференциальное уравнение с произвольным отклонением аргумента


%


\begin{align}


\begin{split}


& ({\cal L}x)(t) \equiv \Dot x(t) + p(t)x[h(t)] = f(t),


\quad t \in [0,T], \\


& x(\xi ) = \phi (\xi ), \quad \xi \notin [0,T];


\end{split}


\end{align}


%


$p,f \in L$, $\phi \in L_\infty$, $h:[0,T] \to R$ --- измеримая функция.


С учетом определения оператора внутренней суперпозиции исходное


уравнение можно записать в виде


%


$$


({\cal L}x)(t) \equiv \dot x(t) + p(t)(S_hx)(t) = f(t) -


p(t)\phi^h(t),  \quad t \in [0,T].


$$


%


В этом случае 


$(Kz)(t) = - \int_0^T p(t) \chi (t,s) z(s)ds$,


где $\chi (t,s)$ --- характеристическая функция множества 


$\{(t,s) \in [0,T] \times [0,T]: 0 \le s \le h(t) \le T \}$.





Отметим, что в примере 1 оператор $Q$ обратим, а в примере 2 оператор $Q$


не является, вообще говоря, обратимым. Обратимость оператора $Q$ имеет место


в случае $h(t) \le t$, $t\in [0,T]$ (случай запаздывающего аргумента).





\section{Основная идея исследования}





Обозначив $z = \Dot x$, запишем уравнение (3) в интегральной форме


%


\begin{equation}


z(t) - \int_0^T k(t,s) z(s) ds = g(t),


\quad t \in [0,T],


\end{equation}


%


где 


$k(t,s)=-p(t)\chi (t,s)$;


$g(t)=f(t)-p(t)\chi (t,0)x(0)-p(t)\phi^h(t)$.


Это уравнение Фредгольма 2-го рода с вполне непрерывным


оператором 


$K:L\to L$; $(Kz)(t) = \int_0^T k(t,s)z(s)ds$


(\cite{Az}, c.\,141). Как известно, ядро $k(t,s)$ может быть с любой наперед


заданной точностью аппроксимировано вырожденным ядром


%


\begin{equation}


k_a(t,s)  = \sum_{i=1}^n u_i (t)  v_i (s),


\quad u_i \in L, \ v_i\in L_\infty, \ i=1,\dotsc,n.


\end{equation}


%


Пусть при этом $\|\Delta K\|_{L \to L} \le \delta_0$, где


%


$$


(\Delta Kz)(t)  =  \int_0^T \{ k(t,s) - k_a(t,s) \} z(s) ds.


$$


%


Как известно (напр., \cite{Zab}), для интегральных уравнений вида


%


\begin{equation}


z(t) - \int_0^T k_a(t,s) z(s) ds = g(t)


\end{equation}


%


однозначная разрешимость (обратимость оператора $[I -K_a]$)


эквивалентна обратимости некоторой конечномерной матрицы $\Lambda$.


Предположим, что обратимость оператора $[I -K_a]$ установлена


и построен резольвентный оператор ${\cal R}_a =[I -K_a]^{-1}-I$:


%


$$


({\cal R}_a g)(t) = \int_0^T r_a(t,s) g(s) ds,


$$


%


дающий представление для решения $z$ уравнения (6),


%


$$


z(t) =g(t) + \int_0^T r_a(t,s)g(s) ds.


$$


%


Пусть $\|I+{\cal R}_a\|_{L \to L} \le \rho_0 $.


Тогда по теореме об обратном операторе


(напр., \cite{Hut}, с.\,99)


в случае, если выполнено неравенство $\delta_0<1/{\rho_0}$,


оператор $[I - K]$ также обратим.





Построение операторов $[I-K_a]$, ${\cal R}_a $, матрицы


$\Lambda$, констант $\delta_0$, $\rho_0$ для дифференциальных уравнений c


произвольным отклонением аргумента и является предметом дальнейшего


рассмотрения.





\section{Описание объекта исследования}





Исследуется задача Коши для дифференциального уравнения с произвольным


отклонением аргумента


%


\begin{align}


\begin{split}


& ({\cal L}x)(t) \equiv \Dot x(t)+p(t)x[h(t)]=f(t),


\quad t \in [0,T], \\


& x(\xi ) =0, \ \xi \notin [0,T]; \ x(0)=\alpha.


\end{split}


\end{align}


%


Здесь $p, f \in L$; $h$ --- кусочно-непрерывная функция, точки


разрыва которой фиксированы, их число конечно, допускаются только разрывы


1-го рода; $\alpha \in R$. Тот факт, что в задаче (7) для упрощения


выкладок взята нулевая начальная функция ($\phi(\xi)\equiv0$), не уменьшает


общности получаемого результата, т.\,к. обратимость главной части


исследуемого дифференциального оператора не зависит от вида начальной


функции. Запишем задачу (7) в интегральной форме


%


\begin{equation}


z(t) - \int_0^T -p(t) \chi (t,s) z(s) ds = g(t),


\quad t \in [0,T].


\end{equation}


%


Здесь $g(t) = f(t) - p(t) \chi (t,0) \alpha$.





\section{Схема исследования разрешимости задачи (7)}





Далее поэтапно описывается конструктивная схема исследования разрешимости


задачи Коши (7).





5.1. {\it Аппроксимация исходной задачи.}


Пусть $\tau_1, \dotsc, \tau_{n_\tau}$ --- точки разрыва функции $h$ и корни


уравнений $h(t)=0$; $h(t)=T$; $t \in [0,T]$; и пусть


$\tau_1^r, \dotsc,\tau_{n_\tau^r}^r$


--- рациональные из этих точек. Определим константу


$\nu=\text{НОД}( T, \tau_1^r, \dotsc, \tau_{n_\tau^r}^r)10^{-\mu}$


$(\mu > 0)$.


Определим $n=T/{\nu}$ и построим систему


точек $t_i = i \cdot \nu$; $i=0, \dotsc, n $; обозначим


$B_i=[t_{i-1}, t_i)$, $i=1,\dotsc, n-1$; 


$B_{n}=[t_{n-1},T]$; $\chi_i(\cdot)$ --- характеристическая


функция множества $B_i$; $i=0, \dotsc, n$.


Параметр $\mu$ выбирается из условия минимизации ошибки


аппроксимации, а также его значение должно быть таким, чтобы каждое из


множеств $B_j$ либо не содержало иррациональных точек $\tau_i$, либо


содержало только одну из них.





Зафиксируем $i=1,\dotsc, n$; $t \in B_i$; функция $h(\cdot)$


аппроксимируется рациональной константой $h_i^a$ с оценкой погрешности


$h_i^v \ge |h(t)-h_i^a|$; функция $p(\cdot)$ аппроксимируется многочленом


с рациональными коэффициентами $p_i^a(\cdot)$ c оценкой погрешности


$p_i^v \ge \int_{t_{i-1}}^{t_i} |p(s)-p_i^a(s)| ds$;


функция $f(\cdot)$ аппроксимируется многочленом


с рациональными коэффициентами $f_i^a(\cdot)$ c оценкой погрешности


$f_i^v\ge \int_{t_{i-1}}^{t_i} |f(s)-f_i^a(s)| ds$; константа $\alpha$


приближается рациональной константой $\alpha_a$ с оценкой погрешности


$\alpha_v \ge |\alpha - \alpha_a|$. Обозначим


%


$$


p_a(t) = \sum_{i=1}^{n} \chi_i(t) p_i^a(t); \quad


f_a(t) = \sum_{i=1}^{n} \chi_i(t) f_i^a(t); \quad


h_a(t) = \sum_{i=1}^{n} \chi_i(t) h_i^a. \quad


$$


%





Запишем аппроксимирующую (7) задачу Коши


%


\begin{align}


\begin{split}


& ({\cal L}_a x)(t) \equiv \Dot x(t) + p_a(t)x[h_a(t)] = f_a(t),


\quad t \in [0,T], \\


& x(\xi ) = 0, \ \xi \notin [0,T]; \ x(0)=\alpha_a;


\end{split}


\end{align}


%


или в интегральной форме


%


\begin{equation}


z(t) - \int_0^T -p_a(t) \chi_a (t,s) z(s) ds = g_a(t),


\quad t \in [0,T],


\end{equation}


%


где $\chi_a (t,s)$ --- характеристическая функция множества


$\{(t,s) \in [0,T] \times [0,T]: \, 0 \le s \le h_a(t) \le T \}$;


$g_a(t) =f_a(t) - p_a(t) \chi_a (t,0) \alpha_a$.


Определим $k_a(t,s) = -p_a(t) \chi_a(t,s)$.


В силу определения $k_a(t,s)$ --- вырожденное ядро вида (5), где


%


\begin{equation}


u_i(t) =-\chi_i(t)  p_i^a(t); \quad


v_i(s) = \chi_{[0, s_{n_i^s}]}(s);


\end{equation}


%


$s_{n_i^s}=


\begin{cases}


h_i^a, & h_i^a \in [0,T]; \\


0, & h_i^a \notin [0,T],


\end{cases}$


$\chi_{[0,\tau]}(\cdot )$ --- характеристическая функция отрезка $[0,\tau]$.





5.2. {\it Решение интегрального уравнения} (10) {\it и построение


резольвентного оператора.}


Следуя известной методике \cite{Zab},


решение $z_a$ уравнения $(10)$ представим


в виде


%


$$


z_a(t) = \sum_{i=1}^{n} u_i(t) \gamma_i + g_a(t).


$$


%


Здесь $n$-вектор 


$\gamma = {\mathrm {col}}\{ \gamma_1, \dotsc,\gamma_{n}\}$


есть решение системы алгебраических уравнений


%


\begin{equation}


\Lambda  \gamma  =  \beta,


\end{equation}


%


где


$\beta = {\mathrm {col}}\{ \beta_1, \dotsc, \beta_{n}\}$;


$\beta_i = \int_0^T v_i(s)g_a(s) ds$,


$i=1, \dotsc, n$;


$\Lambda = \{ \lambda_{ij} \}_{i,j=1}^{n}$;


$\lambda_{ij} = \delta_{ij}-\int_0^T u_j(s) v_i(s) ds$;


$\delta_{ij}$ --- символ Кронекера.





Пусть матрица $\Lambda$ обратима и матрица $M=\Lambda^{-1}$;


$M =\{m_{ij}\}_{i,j=1}^{n}$. Тогда $\gamma_i = \sum_{j=1}^{n}m_{ij} \beta_j$


и


%


$$


z_a(t)=[(I-K_a )^{-1} g_a] (t) =


[(I+{\cal R}_a )  g_a ] (t),


$$


%


где


%


$$


({\cal R}_ag_a ) (t) = \int_0^T r_a(t,s) g_a(s)


ds, \quad r_a(t,s) = -\sum_{i=1}^{n} \sum_{j=1}^{n}


u_i(t) v_j(s) m_{ij}.


$$


%





5.3. {\it Построение оценок норм операторов


$\Delta K _{L \to L}$, $[I+{\cal R}_a]_{L \to L}$.}


Для оценки нормы интегрального оператора


$K:L \to L$ воспользуемся известным соотношением (\cite{Zab}, c.\,107)


%


$$


\| K \| _{L \to L} = \vraimax_s \int_0^T |k(t,s)| dt.


$$


%


Далее имеем


%


\begin{align*}


\| {\cal R}_a \| _{L \to L} &


=  \vraimax_s \int_0^T |r_a(t,s)| dt \le \\


& \le  \vraimax_s \sum_{j=1}^{n} v_j(s)\sum_{i=1}^{n}


\int_{t_{i-1}}^{t_i} |m_{ij} p_i^a (t)| dt \le \\


& \le \vraimax_s \sum_{j=1}^{n} v_j(s) \Tilde m_j,


\end{align*}


%


где рациональная константа $\Tilde m_j$ определяется неравенством


%


$$


\Tilde m_j \ge \sum_{i=1}^{n} |m_{ij}| \sqrt{\nu} \biggl(\,


\int_{t_{i-1}}^{t_i} |p_i^a (t)|^2 dt \biggr) ^{1/2}.


$$


%


Окончательно получаем


%


\begin{equation}


\| I+{\cal R}_a \| _{L \to L} \, \le \, 1+\rho_0, \quad


\rho_0 = \sum_{j=1}^{n} \Tilde m_j.


\end{equation}





Норма оператора $\Delta K: L \to L$,


%


$$


(\Delta Kz) (t) = 


\int_0^T [ p(t)\chi(t,s) - p_a(t)\chi _a(t,s) ] z(s) ds,


$$


%


определится аналогичным образом


%


$$


\| \Delta K \| _{L \to L} = \vraimax_s \int_0^T


 | p(t)\chi(t,s) - p_a(t)\chi _a(t,s) | dt.


$$


%


Поэтому


%


\begin{equation}


\| \Delta K \| _{L \to L}


= \vraimax_s \biggl\{ \int_0^T


 | p_a(t)|\,|\chi(t,s) - \chi _a(t,s) | dt +


\int_0^T  | p(t) - p_a(t) | \chi(t,s) dt \biggr\}.


\end{equation}





Пусть $t \in B_i$, $i=1,\dotsc, n$. Определим константы


%


$$


{\underline d}_i = \max \left\{ 0, h_i^a -h_i^v \right\}; \quad


\bar d_i = \min \left\{ T, h_i^a +h_i^v \right\}.


$$


%


Справедливы оценки


%


$$


|\chi (t,s) - \chi_a(t,s) | \le \chi _{[{\underline d}_i, \bar d_i]} (s)


 \chi _i(t); \quad


\chi (t,s) \le \chi _{[0, \bar d_i]} (s)  \chi _i(t).


$$





С учетом последних неравенств продолжим оценку (14)


%


$$


\| \Delta K \| _{L \to L}


\le \vraimax_s \biggl\{ \sum_{i=1}^n


\int_{t_{i-1}}^{t_i}


| p_a(t)| dt\, \chi _{[{\underline d}_i, \bar d_i]}(s) +


\sum_{i=1}^n


\int_{t_{i-1}}^{t_i} | p(t) - p_a(t) | dt\,


\chi _{[0, \bar d_i]}(s)\biggr\} \le \delta_0,


$$


%


где рациональная константа $\delta_0$ определяется


неравенством


%


\begin{equation}


\delta_0 \ge \sum_{i=1}^n p^v_i + \max_j


\bigl\{ p^N_{q^j_1}+\dotsb+p^N_{q^j_{n_j}} \bigr\}.


\end{equation}


%


Здесь


$p_{\tau}^N \ge \sqrt{\nu} \Bigl(


\int_{t_{\tau-1}}^{t_{\tau}} |p_i^a (t)|^2 dt \Bigr)^{1/2}$,


а максимум ищется по всем $j$, для которых


$\mathrm {mes}\,S_j\ne 0$,


где $S_j=


\left\{ [ \underline d_{q^j_1}, \bar d_{q^j_1}] \cap \dotsc


\cap [ \underline d_{q^j_{n_j}}, \bar d_{q^j_{n_j}}]


\right\}$.





5.4. {\it Проверка условия разрешимости. Построение оценки погрешности


решения.} Выполнение условия


%


\begin{equation}


\delta_0 < \frac {1}{1+\rho_0}


\end{equation}


%


обеспечивает однозначную всюду разрешимость исследуемой задачи


$(7)$. Кроме того, имеют место оценки


%


\begin{align}


\begin{split}


\| I+{\cal R} \|_{L\to L} & \le 1 + \rho_0 + \frac {\delta_0


(1+\rho_0)^2} {1-\delta_0 (1+\rho_0)}; \\


%


\| z -z_a\|_L & \le \frac {\delta_0 (1+\rho_0)^2}


{1 - \delta_0 (1+\rho_0 )}


(g^N + g^v) + (1+\rho_0)g^v,


\end{split}


\end{align}


%


где $[I+{\cal R}] = [I-K]^{-1}$; $z$ ---


решение уравнения $(8)$; константы $g^N$ и $g^v$ определяются неравенствами


%


$$


g^N \ge \int_0^T |g_a(t)| dt, \qquad


g^v \ge \int_0^T |g(t) - g_a(t)| dt.


$$





\begin{notenonum} Пусть задача (7) однозначно всюду разрешима. Тогда в


силу сделанных предположений о параметрах задачи (7) всегда можно


построить систему функций вида (11) такую, что матрица $\Lambda$ системы


(12) будет обратима, а константы $\delta_0$ и $\rho_0$, определяемые


соотношениями (15) и (13) соответственно, будут удовлетворять


неравенству (16).


\end{notenonum}





\section{Иллюстрирующий пример}





Предложенный выше метод исследования однозначной разрешимости задачи вида


(7) был реализован средствами системы MAPLE. С помощью разработанной


программы, в частности, был проведен анализ модельного примера


%


\begin{align}


\begin{split}


& \dot x(t) + p(t)x[h(t)] = t^2,


\quad t \in [0,1], \\


& x(\xi ) = 0, \ \xi \notin [0,1];


\end{split}


\end{align}


%


$$


p(t) =


\begin{cases}


2t, &  t \in [0,0.3); \\


t/4, & t \in [0.3,0.7); \\


2t, & t \in [0.7,1],


\end{cases}


\qquad


%


h(t) = 


\begin{cases}


0.65, & t \in [0,0.3); \\


1-t, & t \in [0.3,0.7); \\


0.8, & t \in [0.7,1].


\end{cases}


$$


%


Было достоверно показано, что задача (18) однозначно всюду


разрешима, при этом были получены оценки


%


$$


\delta_0 \le 0.0008; \quad \rho_0 \le 0.8025; \quad


\frac {1}{1+\rho_0} \ge 0.5547;


$$


%


$$


\|I+{\cal R}\|_{L \to L} \le 1.805; \quad


\|z - z_a\|_L \le 0.009; \quad \|z\|_L \le 0.7831.


$$








\begin{thebibliography}{99} 


 


\bibitem{Az}


Азбелев Н.В., Максимов В.П., Рахматуллина Л.Ф. {\em Введение в теорию


функционально-дифференциальных уравнений}. -- М.: Наука, 1991. -- 280 с.





\bibitem{Mak1}


Максимов В.П., Румянцев А.Н. {\em Краевые задачи и задачи импульсного


управления в экономической динамике. Конструктивное исследование} //


Изв. вузов. Математика. -- 1993. -- \No\,5. -- С.\,56--71.





\bibitem{Ru1}


Румянцев А.Н. {\em Исследование разрешимости краевых задач с применением


векторных априорных неравенств при наличии импульсных возмущений} //


Изв. вузов. Математика. -- 1994. -- \No\,6. -- С.\,63--77.





\bibitem{Ru2}


Румянцев А.Н., Федяев Е.А. {\em Конструктивное исследование разрешимости


краевых задач для систем дифференциальных уравнений с сосредоточенным


запаздыванием при наличии импульсных воздействий} // Изв. вузов.


Математика. -- 1995. -- \No\,10. -- С.\,46--60.





\bibitem{Mak2}


Maksimov V., Rumyantsev A. {\em The study of the solvability of BVP's for


FDE's: A constructive approach} // Boundary Value Problems for Functional


Differential Equations. World Scientific Publ., Singapore, 


1995. -- P.\,227--237.





\bibitem{Kau1}


Kaucher E.W., Miranker W.L. {\em Self-validating numerics for function


space problems}. -- New York: Academic Press, 1984.





\bibitem{Kau2}


Kaucher E.W., Miranker W.L. {\em Validating computation in a function


space.} // Reability in computing. -- New York: Academic Press, 


1988. -- P.\,403--425.





\bibitem{Moo}


Moore R.E. {\em Shen Zuhe interval methods for operator equations} //


Reability in computing. -- New York: Academic Press, 1988. -- P.\,379--389.





\bibitem{Plum}


Plum M. {\em Computer-assisted existence proofs for two-point boundary


value problems}. -- Computing. Springer-Verlag, 1991.





\bibitem{Zab}


Забрейко П.П. {\em Интегральные уравнения}. -- М.: 1968. -- 448 с.





\bibitem{Hut}


Хатсон~В., Пим~Дж. {\em Приложения функционального анализа и теории


операторов}. -- М.: Мир, 1983. -- 431~с.





\end{thebibliography} 


 


\vskip 2cm 


 


\post{\it Пермский государственный}{\it Поступила}


\post{\it университет}{28.06.1996} 


 


%% use \post in Eng. 


\end{document} 


